\chapter{Inequalities in Two Variables: Regions in the Plane}

The algebra volume solved inequalities in one variable, where the answer is
a set of points on a line --- an interval, or a union of a few. With two variables
the same question has a two-dimensional answer: the solution set of
$2x+3y\le12$ is not a list of numbers but a \emph{region} of the plane, and
it is described by drawing it.

The mechanism is short. The corresponding \emph{equation} $2x+3y=12$ is a
curve, and a curve cuts the plane into pieces. On each piece the expression
$2x+3y-12$ cannot change sign without passing through zero, so its sign is
constant there. One test point per piece therefore settles the whole
region --- the same sign-table logic as in one variable, with curves in
place of points.

Two directions are worth drilling, and examinations favour the second. The
forward direction is: given inequalities, shade the region. The reverse is:
given a shaded region, write down the inequalities that describe it. The
reverse is harder and more useful, because it is what modelling actually
requires --- a feasible set is described in words or in a picture long before
anyone writes the constraints. It is also exactly the skill that the
graphical solution of a linear programme depends on.

\section{Core Rules}

\begin{rulebox}{The solution set is a region}
For an expression $F(x,y)$, the set $\{F(x,y)<0\}$ is bounded by the curve
$\{F(x,y)=0\}$. The recipe:
\begin{enumerate}[leftmargin=1.8em,itemsep=0.3ex]
  \item \textbf{Draw the boundary} $F(x,y)=0$. It is one of the curves of
        Chapter~2: a line, circle, ellipse, parabola or hyperbola.
  \item \textbf{Note which pieces it creates.} A line makes two half-planes;
        a circle makes an inside and an outside; a parabola makes an
        ``above'' and a ``below''.
  \item \textbf{Test one point per piece.} Substituting a single convenient
        point --- the origin whenever it is not on the boundary --- decides
        the sign for that whole piece.
  \item \textbf{Mark the boundary.} Solid if the inequality is
        $\le$ or $\ge$ (the boundary belongs to the set), dashed if it is
        $<$ or $>$ (it does not).
\end{enumerate}
\end{rulebox}

\begin{rulebox}{The standard regions}
\begin{itemize}[leftmargin=1.6em,itemsep=0.3ex]
  \item \textbf{Half-plane:} $ax+by\le c$. Draw the line, test the origin.
        With $c>0$ and the origin satisfying it, the region is the side
        containing the origin.
  \item \textbf{Disc and its complement:} $x^{2}+y^{2}\le r^{2}$ is the
        closed disc; $\ge$ gives everything outside. Shifted:
        $(x-x_{0})^{2}+(y-y_{0})^{2}\le r^{2}$.
  \item \textbf{Ellipse interior:} $\dfrac{x^{2}}{a^{2}}
        +\dfrac{y^{2}}{b^{2}}\le1$.
  \item \textbf{Above or below a parabola:} $y\ge x^{2}$ is the region
        \emph{inside} the parabola's opening; $y\le x^{2}$ is everything
        outside it. Rewriting an inequality as $y\ge\cdots$ or
        $y\le\cdots$ makes the side unambiguous: ``$y$ at least this'' means
        upwards.
  \item \textbf{Between hyperbola branches:} $xy\ge c$ with $c>0$ is the
        pair of regions beyond the two branches --- and on the positive
        quadrant alone it is the convex ``better set'' of Cobb--Douglas.
  \item \textbf{Strip:} $|x|\le a$ is the vertical strip $-a\le x\le a$;
        $|y|\le b$ a horizontal one.
\end{itemize}
\end{rulebox}

\begin{gallery}
\begin{tikzpicture}
  \begin{axis}[geoaxis, xmin=-1.4, xmax=7.4, ymin=-1.4, ymax=5.4,
      xtick={2,4,6}, ytick={2,4}, title={$2x+3y\le12$, $x,y\ge0$}]
    \addplot[draw=none, name path=lin, domain=0:6, samples=2]{(12-2*x)/3};
    \addplot[draw=none, name path=ax, domain=0:6, samples=2]{0};
    \region{lin}{ax}{}
    \addplot[domain=-1.2:7.2, samples=2]{(12-2*x)/3};
  \end{axis}
\end{tikzpicture}\hfill
\begin{tikzpicture}
  \begin{axis}[geoaxis, xmin=-3.4, xmax=3.4, ymin=-3.4, ymax=3.4,
      xtick={-2,2}, ytick={-2,2}, title={$x^{2}+y^{2}\le4$}]
    \addplot[draw=none, name path=up, domain=-2:2, samples=201]{sqrt(4-x^2)};
    \addplot[draw=none, name path=lo, domain=-2:2, samples=201]{-sqrt(4-x^2)};
    \region{up}{lo}{}
    \addplot[domain=0:360, samples=181]({2*cos(x)},{2*sin(x)});
  \end{axis}
\end{tikzpicture}\hfill
\begin{tikzpicture}
  \begin{axis}[geoaxis, xmin=-2.6, xmax=2.6, ymin=-0.8, ymax=4.6,
      xtick={-2,2}, ytick={2,4}, title={$x^{2}\le y\le 4$}]
    \addplot[draw=none, name path=par, domain=-2:2, samples=201]{x^2};
    \addplot[draw=none, name path=top, domain=-2:2, samples=2]{4};
    \region{par}{top}{}
    \addplot[domain=-2.15:2.15, samples=201]{x^2};
    \addplot[domain=-2.5:2.5, samples=2]{4};
  \end{axis}
\end{tikzpicture}
\gcap{Three regions: an intersection of three half-planes (a triangle), a
disc, and the region caught between a parabola and a line.}
\end{gallery}

\begin{rulebox}{Systems: intersect}
Several inequalities holding \emph{simultaneously} give the
\textbf{intersection} of their regions --- shade each one and keep only the
overlap. Two consequences:
\begin{itemize}[leftmargin=1.6em,itemsep=0.3ex]
  \item An intersection of half-planes is a \textbf{convex polygon}
        (possibly unbounded). This is precisely the feasible region of a
        linear programme, and the reason the optimum of an LP sits at a
        corner.
  \item An intersection of convex regions is always convex, so discs,
        ellipse interiors and half-planes may be combined freely and the
        result stays convex. A \emph{union}, by contrast, usually is not,
        and neither is the outside of a disc.
\end{itemize}
\end{rulebox}

\begin{rulebox}{The reverse direction: from picture to inequalities}
Given a shaded region, work boundary by boundary.
\begin{enumerate}[leftmargin=1.8em,itemsep=0.3ex]
  \item \textbf{Identify every boundary curve} and write its equation. A
        circle needs its centre and radius; a line two points or an
        intercept and a slope; a parabola its vertex and one further
        point.
  \item \textbf{Fix the direction of each} with a point that is clearly
        inside the shaded region: substitute it and keep the inequality
        sign that comes out true.
  \item \textbf{Join the conditions with ``and''}, one per boundary. A region
        bounded by three curves needs three inequalities.
  \item \textbf{Check a point just outside} each boundary: it must fail
        exactly one condition. If it fails none, a condition is missing.
\end{enumerate}
Read the boundary style too: a dashed boundary means a strict inequality.
\end{rulebox}

\begin{rulebox}{Which of these regions are convex}
One rule settles almost every case. For a function $f$ of one variable:
\begin{itemize}[leftmargin=1.6em,itemsep=0.3ex]
  \item $f$ \textbf{convex} $\Longrightarrow$ the region \emph{above} it,
        $\{y\ge f(x)\}$ (its \de{Epigraph}), is a convex set. The region
        \emph{below} it is then not convex.
  \item $f$ \textbf{concave} $\Longrightarrow$ the region \emph{below} it,
        $\{y\le f(x)\}$, is convex. The region above it is then not.
\end{itemize}
The proof is one line: if $y_{i}\ge f(x_{i})$ and $f$ is convex then
$ty_{1}+(1-t)y_{2}\ge tf(x_{1})+(1-t)f(x_{2})\ge f(tx_{1}+(1-t)x_{2})$.
So $y\ge x^{2}$ is convex and $y\le x^{2}$ is not, while $y\le4-x^{2}$
\emph{is} convex, because $4-x^{2}$ is concave. Add to this that
half-planes, discs and ellipse interiors are convex, that any
\textbf{intersection} of convex sets is convex, and that the
\emph{outside} of a disc and any \emph{union} of separated pieces is not,
and every region in this chapter is decided.
\end{rulebox}

\begin{rulebox}{Absolute values in two variables}
Split on the sign, exactly as in one variable, or recognise the standard
shapes:
\[
  |x|+|y|\le c \;\text{(a diamond, or square rotated }45^{\circ}),\qquad
  \max(|x|,|y|)\le c \;\text{(a square)} .
\]
For $|y|\le f(x)$ the region is squeezed between $y=f(x)$ and $y=-f(x)$,
and exists only where $f(x)\ge0$. For $y\ge|x|$ it is the wedge above both
lines $y=x$ and $y=-x$.
\end{rulebox}

\section{Worked Example}

\begin{worked}{Sketch $S=\{(x,y):x^{2}+y^{2}\le16,\ y\ge0,\ x\le2\}$, then decide whether $\{(x,y):y\ge\sqrt{x}\}$ is convex}
\wstep{Step 1. Take the constraints one at a time, and shade each.}
\begin{itemize}[leftmargin=1.6em,itemsep=0.2ex]
  \item $x^{2}+y^{2}\le16$: the closed disc of radius $4$ about the origin
        --- ``$\le$'' is the inside, because the origin satisfies it
        ($0\le16$).
  \item $y\ge0$: the upper half-plane, boundary included.
  \item $x\le2$: everything to the left of the vertical line $x=2$.
\end{itemize}
\wtrap{Decide inside-versus-outside by testing one point, not by staring at
the inequality sign. The origin is free and settles the disc immediately.
For a line not through the origin, any convenient point does.}

\wstep{Step 2. Intersect the three.}
The region is the part of the upper half-disc lying left of $x=2$: a
half-disc with a circular segment removed on the right. Its boundary
consists of three pieces --- the segment of the $x$-axis from $x=-4$ to
$x=2$, the vertical segment $x=2$ from $y=0$ up to $y=\sqrt{16-4}=2\sqrt3$,
and the arc of the circle from $(2,2\sqrt3)$ counter-clockwise round to
$(-4,0)$.

\wstep{Step 3. Find the corner points explicitly.}
Where $x=2$ meets the circle: $4+y^{2}=16$, so $y=2\sqrt3\approx3.46$ (the
negative root is cut off by $y\ge0$). Where $y=0$ meets the circle:
$x=\pm4$, of which $x=4$ is cut off by $x\le2$, leaving $(-4,0)$. And
$(2,0)$ where the two straight boundaries meet. Three corners:
$(-4,0)$, $(2,0)$, $(2,2\sqrt3)$.

\wstep{Step 4. Sketch, and check one interior point.}
\begin{center}
\begin{tikzpicture}
\begin{axis}[geoaxis, width=6.6cm, height=6.6cm,
             xmin=-5.2, xmax=5.2, ymin=-1.2, ymax=5.2,
             xtick={-4,-2,2,4}, ytick={2,3.464},
             yticklabels={$2$,$2\sqrt3$}]
  % the region: under the upper semicircle, above y=0, left of x=2
  \addplot[accent, very thick, domain=-4:2, samples=200,
           fill=regionfill, fill opacity=0.5] {sqrt(16-x^2)} \closedcycle;
  % the discarded part of the disc, for contrast
  \addplot[accent!35, thick, dashed, domain=2:4, samples=80]
           {sqrt(16-x^2)};
  \node[font=\scriptsize] at (axis cs:-1.2,1.4) {$S$};
  \node[font=\tiny, anchor=south west] at (axis cs:2.1,3.6) {$(2,2\sqrt3)$};
\end{axis}
\end{tikzpicture}
\end{center}
Test $(0,1)$: $0+1\le16$ \checkmark, $1\ge0$ \checkmark, $0\le2$
\checkmark. Inside, as drawn.

\wstep{Step 5. Note that $S$ is convex, and why you can say so instantly.}
Each of the three constraint sets is convex --- a disc, a half-plane, a
half-plane --- and \textbf{an intersection of convex sets is convex}. No
picture is needed for that conclusion, and no amount of picture-gazing
substitutes for it.

\wstep{Step 6. The second region: $\{y\ge\sqrt{x}\}$. Do not trust the eye.}
The boundary $y=\sqrt{x}$ is a smooth, gently bending curve, and the region
above it looks perfectly well behaved. Test the definition instead. Take the
two boundary points
\[
  P=(0,0),\qquad Q=(4,2),
\]
both in the set, since $0\ge\sqrt0$ and $2\ge\sqrt4$. Their midpoint is
$M=(2,1)$, and
\[
  \sqrt{2}\approx1.414>1 ,
\]
so $M$ violates $y\ge\sqrt{x}$ and lies \emph{outside}. A segment between
two points of the set leaves the set:
\[
  \boxed{\{(x,y):y\ge\sqrt{x}\}\ \text{is not convex.}}
\]

\wstep{Step 7. The principle behind it.}
$\sqrt{x}$ is \textbf{concave}. The set above the graph of a function --- its
epigraph --- is convex exactly when the function is \emph{convex}; the set
below --- the hypograph --- is convex exactly when the function is
\emph{concave}. So $\{y\ge\sqrt{x}\}$ fails while $\{y\le\sqrt{x}\}$
succeeds, and $\{y\ge x^{2}\}$ succeeds while $\{y\le x^{2}\}$ fails.
\wtrap{``Above'' is not a synonym for ``convex''. Match the direction of the
inequality to the curvature of the function: convex function with $\ge$, or
concave function with $\le$. Getting this backwards is the single most
common mistake in the chapter, and it recurs in Volume~7 where the feasible
set of a constrained problem must be convex for the first-order conditions
to be sufficient.}
\end{worked}

\section{Practice Problems}

Sketch the region described by each inequality in problems 1--12, and say
whether it is convex.

\begin{problems}
  \item $y\ge 2x-1$
  \item $x+y\le 3$
  \item $x^{2}+y^{2}\le 9$
  \item $x^{2}+y^{2}>4$
  \item $y\ge x^{2}$
  \item $y\le 4-x^{2}$
  \item $\dfrac{x^{2}}{9}+\dfrac{y^{2}}{4}\le1$
  \item $xy\ge 1$ on $x>0$
  \item $|x|\le 2$
  \item $|y|\le x$
  \item $(x-2)^{2}+(y+1)^{2}\le 9$
  \item $y>\dfrac1x$ on $x>0$
\end{problems}

In problems 13--20 sketch the region defined by the whole system.

\begin{problems}[resume]
  \item $x\ge0$, \;$y\ge0$, \;$2x+3y\le12$
  \item $x\ge0$, \;$y\ge0$, \;$x+y\le4$, \;$y\le2$
  \item $x^{2}+y^{2}\le16$, \;$y\ge0$
  \item $x^{2}+y^{2}\le16$, \;$y\ge0$, \;$x\le2$
  \item $x^{2}\le y$, \;$y\le x+2$
  \item $1\le x^{2}+y^{2}\le9$
  \item $x^{2}+y^{2}\le9$, \;$x\ge0$, \;$y\ge0$
  \item $y\ge x^{2}$, \;$y\le 4$, \;$x\ge0$
\end{problems}

Problems 21--28 refer to the eight shaded regions below. For each, give a
system of inequalities describing it.

\begin{gallery}
\begin{tikzpicture}
  \begin{axis}[geoaxis, xmin=-5.2, xmax=5.2, ymin=-5.2, ymax=5.2,
      xtick={-4,-2,2,4}, ytick={-4,-2,2,4}, title={(a)}]
    \addplot[draw=none, name path=up, domain=-4:2, samples=201]{sqrt(16-x^2)};
    \addplot[draw=none, name path=ax, domain=-4:2, samples=2]{0};
    \region{up}{ax}{}
    \addplot[domain=0:360, samples=181]({4*cos(x)},{4*sin(x)});
    \addplot[gray!60, dashed, thin, samples=2] coordinates {(2,-5)(2,5)};
  \end{axis}
\end{tikzpicture}\hfill
\begin{tikzpicture}
  \begin{axis}[geoaxis, xmin=-2.8, xmax=2.8, ymin=-0.9, ymax=5.1,
      xtick={-2,2}, ytick={2,4}, title={(b)}]
    \addplot[draw=none, name path=par, domain=-2:2, samples=201]{x^2};
    \addplot[draw=none, name path=top, domain=-2:2, samples=2]{4};
    \region{par}{top}{}
    \addplot[domain=-2.2:2.2, samples=201]{x^2};
    \addplot[domain=-2.6:2.6, samples=2]{4};
  \end{axis}
\end{tikzpicture}\hfill
\begin{tikzpicture}
  \begin{axis}[geoaxis, xmin=-3.9, xmax=3.9, ymin=-2.9, ymax=2.9,
      xtick={-2,2}, ytick={-2,2}, title={(c)}]
    \addplot[draw=none, name path=ell, domain=0:3, samples=201]{2*sqrt(1-x^2/9)};
    \addplot[draw=none, name path=ax, domain=0:3, samples=2]{0};
    \region{ell}{ax}{}
    \addplot[domain=0:360, samples=181]({3*cos(x)},{2*sin(x)});
  \end{axis}
\end{tikzpicture}
\end{gallery}

\begin{gallery}
\begin{tikzpicture}
  \begin{axis}[geoaxis, xmin=-2.4, xmax=3.4, ymin=-0.9, ymax=4.9,
      xtick={-2,2}, ytick={2,4}, title={(d)}]
    \addplot[draw=none, name path=par, domain=-1:2, samples=201]{x^2};
    \addplot[draw=none, name path=lin, domain=-1:2, samples=2]{x+2};
    \region{par}{lin}{}
    \addplot[domain=-2.1:2.2, samples=201]{x^2};
    \addplot[domain=-2.2:3.2, samples=2]{x+2};
  \end{axis}
\end{tikzpicture}\hfill
\begin{tikzpicture}
  \begin{axis}[geoaxis, xmin=-3.9, xmax=3.9, ymin=-3.9, ymax=3.9,
      xtick={-2,2}, ytick={-2,2}, title={(e)}]
    \addplot[draw=none, name path=o1, domain=-3:3, samples=201]{sqrt(9-x^2)};
    \addplot[draw=none, name path=o2, domain=-3:3, samples=201]{-sqrt(9-x^2)};
    \region{o1}{o2}{}
    \addplot[draw=none, name path=i1, domain=-1:1, samples=201]{sqrt(1-x^2)};
    \addplot[draw=none, name path=i2, domain=-1:1, samples=201]{-sqrt(1-x^2)};
    \addplot[white, forget plot, sharp plot] fill between[of=i1 and i2];
    \addplot[domain=0:360, samples=181]({3*cos(x)},{3*sin(x)});
    \addplot[domain=0:360, samples=181]({cos(x)},{sin(x)});
  \end{axis}
\end{tikzpicture}\hfill
\begin{tikzpicture}
  \begin{axis}[geoaxis, xmin=-1.4, xmax=7.4, ymin=-1.4, ymax=5.4,
      xtick={2,4,6}, ytick={2,4}, title={(f)}]
    \addplot[draw=none, name path=lin, domain=0:6, samples=2]{(12-2*x)/3};
    \addplot[draw=none, name path=ax, domain=0:6, samples=2]{0};
    \region{lin}{ax}{}
    \addplot[domain=-1.2:7.2, samples=2]{(12-2*x)/3};
  \end{axis}
\end{tikzpicture}
\end{gallery}

\begin{gallery}
\begin{tikzpicture}
  \begin{axis}[geoaxis, xmin=-1.2, xmax=5.4, ymin=-2.9, ymax=2.9,
      xtick={2,4}, ytick={-2,2}, title={(g)}]
    \addplot[draw=none, name path=up, domain=0:4, samples=201]{sqrt(x)};
    \addplot[draw=none, name path=lo, domain=0:4, samples=201]{-sqrt(x)};
    \region{up}{lo}{}
    \addplot[domain=0:4.6, samples=201]{sqrt(x)};
    \addplot[domain=0:4.6, samples=201]{-sqrt(x)};
    \addplot[gray!60, dashed, thin, samples=2] coordinates {(4,-2.8)(4,2.8)};
  \end{axis}
\end{tikzpicture}\hfill
\begin{tikzpicture}
  \begin{axis}[geoaxis, xmin=-3.4, xmax=3.4, ymin=-3.4, ymax=3.4,
      xtick={-2,2}, ytick={-2,2}, title={(h)}]
    \addplot[draw=none, name path=up, domain=-2:2, samples=101]{2-abs(x)};
    \addplot[draw=none, name path=lo, domain=-2:2, samples=101]{abs(x)-2};
    \region{up}{lo}{}
    \addplot[domain=-2:2, samples=101]{2-abs(x)};
    \addplot[domain=-2:2, samples=101]{abs(x)-2};
  \end{axis}
\end{tikzpicture}
\gcap{Regions (a)--(h) for problems 21--28. A dashed boundary means the
inequality there is strict.}
\end{gallery}

\begin{problems}[resume]
  \item Region (a).
  \item Region (b).
  \item Region (c).
  \item Region (d).
  \item Region (e).
  \item Region (f).
  \item Region (g).
  \item Region (h).
\end{problems}

\begin{problems}[resume]
  \item Which of the eight regions (a)--(h) are convex?
  \item Sketch $|x|+|y|\le 3$ and give its four vertices.
  \item Sketch $y\ge|x|$ and describe it in words.
  \item Sketch $|y|\le 4-x^{2}$ and give the $x$-values for which the
        region is non-empty.
  \item Write inequalities for the closed square with vertices
        $(\pm2,\pm2)$, using absolute values and then without them.
  \item A test point satisfies every inequality of a system but lies
        outside the intended region. What has gone wrong?
  \item Explain why the outside of a disc is never convex.
  \item Give a system of two inequalities whose solution set is empty, both
        of them satisfiable on their own.
  \item Show that the intersection of any two convex regions is convex.
  \item Is the union of the two half-planes $y\ge1$ and $y\le-1$ convex?
        Justify.
  \item Write inequalities for the set of points strictly inside the circle
        of radius $2$ about the origin and strictly above the line $y=x$.
  \item For which values of $c$ is the region $\{x^{2}+y^{2}\le c\}$ empty,
        a single point, and a disc?
\end{problems}

\begin{problems}[resume]
  \item \textbf{Budget set.} A consumer has income $60$ and faces prices
        $p_{x}=3$, $p_{y}=4$. Write the budget set as a system, sketch it,
        and give its three vertices.
  \item For the same consumer, write the set of bundles that are
        \emph{affordable and} give at least utility $12$ under
        $U(x,y)=xy$. Sketch it and say whether it is convex.
  \item \textbf{Production.} A firm can use at most $40$ hours of labour
        and $30$ units of capital; product A needs $2$ hours and $1$ unit,
        product B needs $1$ hour and $2$ units. Write the feasible set and
        sketch it.
  \item For the feasible set of the previous problem, give all its
        vertices, and say why a linear objective attains its maximum at one
        of them.
  \item \textbf{Tax wedge.} The set $\{(q,p): p\le 100-2q,\ p\ge 20+q,\
        q\ge0\}$ describes the quantities at which demand price exceeds
        supply price. Sketch it and find the largest $q$ in it.
\end{problems}

\section{Solutions}

\begin{solutions}
  \item The closed half-plane on or above the line $y=2x-1$. Testing
        $(0,0)$: $0\ge-1$ true, so the origin is in it. Convex.
  \item The closed half-plane on or below $x+y=3$ (through $(3,0)$ and
        $(0,3)$); the origin satisfies it. Convex.
  \item The closed disc of radius $3$ about the origin. Convex.
  \item Everything strictly outside the circle of radius $2$, boundary
        dashed and excluded. \textbf{Not} convex.
  \item The region on or above the parabola $y=x^{2}$ --- inside its
        opening. Convex.
  \item The region on or below the downward parabola $y=4-x^{2}$, which
        includes everything far below it. Since $4-x^{2}$ is
        \emph{concave}, the region below it is \textbf{convex}.
  \item The closed elliptical disc with semi-axes $3$ and $2$. Convex.
  \item The region on or beyond the hyperbola branch $y=1/x$ in the first
        quadrant. Convex (this is the quasiconcavity of $xy$).
  \item The closed vertical strip $-2\le x\le2$. Convex, being the
        intersection of two half-planes.
  \item Requires $x\ge0$ and $-x\le y\le x$: the wedge opening to the right
        between the lines $y=\pm x$. Convex.
  \item The closed disc of radius $3$ centred at $(2,-1)$. Convex.
  \item The region strictly above the hyperbola branch $y=1/x$ for $x>0$,
        boundary dashed. Convex.
  \item The triangle with vertices $(0,0)$, $(6,0)$, $(0,4)$ --- the
        intersection of three half-planes.
  \item The quadrilateral with vertices $(0,0)$, $(4,0)$, $(2,2)$,
        $(0,2)$: the triangle $x+y\le4$ in the first quadrant, with its top
        cut off by $y\le2$.
  \item The upper half-disc of radius $4$: the semicircular region above
        the $x$-axis.
  \item The same half-disc with everything right of $x=2$ removed --- the
        region shown in figure (a).
  \item The lens between the parabola and the line, from $x=-1$ to $x=2$
        (the boundaries meet where $x^{2}=x+2$, i.e. $x=-1,2$). Convex.
  \item The annulus between the circles of radius $1$ and $3$. Not convex.
  \item The quarter-disc of radius $3$ in the first quadrant. Convex.
  \item The region above the parabola, below $y=4$, right of the $y$-axis:
        the right half of figure (b).
  \item $x^{2}+y^{2}\le16$, \;$y\ge0$, \;$x<2$. The circle and the
        $x$-axis are solid, so those two are non-strict; the vertical
        boundary is dashed, so that one is strict.
  \item $y\ge x^{2}$ and $y\le4$. (Equivalently
        $x^{2}\le y\le 4$, which also forces $-2\le x\le2$ without needing
        a separate condition.)
  \item $\dfrac{x^{2}}{9}+\dfrac{y^{2}}{4}\le1$, \;$x\ge0$, \;$y\ge0$: the
        quarter of the elliptical disc in the first quadrant.
  \item $y\ge x^{2}$ and $y\le x+2$.
  \item $1\le x^{2}+y^{2}\le9$, that is $x^{2}+y^{2}\ge1$ \emph{and}
        $x^{2}+y^{2}\le9$.
  \item $2x+3y\le12$, \;$x\ge0$, \;$y\ge0$.
  \item $y^{2}\le x$ and $x<4$; equivalently
        $-\sqrt{x}\le y\le\sqrt{x}$ with $x<4$. The parabola opens to the
        right and is solid; the line $x=4$ is dashed.
  \item $|x|+|y|\le2$. Written without absolute values it is the four
        conditions $x+y\le2$, $x-y\le2$, $-x+y\le2$, $-x-y\le2$.
  \item All except (e). Each of (a), (b), (c), (d), (f), (g), (h) is an
        intersection of convex sets: discs and ellipse interiors,
        half-planes, the region above the convex $x^{2}$, the region
        \emph{right} of the parabola $x=y^{2}$ (the same rule with the axes
        exchanged, since $y\mapsto y^{2}$ is convex), and the diamond
        $|x|+|y|\le2$, which is an intersection of four half-planes.
        Only (e), the annulus, fails: it is the disc of radius $3$ with a
        hole, and the segment from $(-2,0)$ to $(2,0)$ runs straight
        through the hole.
  \item A diamond with vertices $(3,0)$, $(0,3)$, $(-3,0)$, $(0,-3)$: the
        square of side $3\sqrt2$ rotated by $45^{\circ}$.
  \item The wedge above both lines $y=x$ and $y=-x$ --- everything at or
        above the V of $y=|x|$, opening upwards with its vertex at the
        origin. Convex.
  \item Squeezed between $y=4-x^{2}$ and $y=-(4-x^{2})$, so it exists only
        where $4-x^{2}\ge0$, i.e. for $-2\le x\le2$. It is the
        lens-shaped region symmetric about both axes.
  \item With absolute values: $\max(|x|,|y|)\le2$, i.e. $|x|\le2$ and
        $|y|\le2$. Without: $-2\le x\le2$ and $-2\le y\le2$, four
        half-planes.
  \item The system describes a \emph{larger} set than intended --- a
        condition is missing. Every boundary of the intended region must
        contribute one inequality; a point outside that satisfies them all
        proves one boundary has been forgotten.
  \item Take two points diametrically opposite each other just outside the
        circle: the segment between them passes through the centre of the
        disc, which is not in the set. So the segment leaves the set and
        convexity fails.
  \item For instance $y\ge x+1$ and $y\le x-1$. Each is a satisfiable
        half-plane, but together they demand $x+1\le y\le x-1$, impossible
        since $x+1>x-1$. The two half-planes are on opposite sides of
        parallel lines.
  \item Let $C_{1},C_{2}$ be convex and take $P,Q\in C_{1}\cap C_{2}$.
        Since $P,Q\in C_{1}$ and $C_{1}$ is convex, the whole segment
        $[P,Q]$ lies in $C_{1}$; the same argument puts it in $C_{2}$.
        Hence $[P,Q]\subseteq C_{1}\cap C_{2}$, which is the definition of
        convexity. (The argument works for any number of sets.)
  \item No. The points $(0,1)$ and $(0,-1)$ both lie in the union, but the
        midpoint $(0,0)$ lies in neither half-plane. A union of disjoint
        convex sets is essentially never convex.
  \item $x^{2}+y^{2}<4$ \;and\; $y>x$. Both boundaries are dashed: the open
        half-disc lying above the diagonal. (It is convex, being an
        intersection of two convex sets.)
  \item Empty for $c<0$; the single point $(0,0)$ for $c=0$; the closed disc
        of radius $\sqrt{c}$ for $c>0$. The three cases are exactly the
        three things a level \emph{set} of $x^{2}+y^{2}$ can do, and they
        are why a sublevel set can be empty without the inequality being
        malformed.
  \item $3x+4y\le60$, \;$x\ge0$, \;$y\ge0$: the triangle with vertices
        $(0,0)$, $(20,0)$ and $(0,15)$.
  \item $3x+4y\le60$, \;$xy\ge12$, \;$x,y>0$: the part of the budget
        triangle lying beyond the indifference curve $y=12/x$. It is the
        intersection of a convex set (the triangle) with another convex set
        (the better set of a quasiconcave utility), hence \textbf{convex}.
  \item With $a$ units of A and $b$ of B: $2a+b\le40$ (labour),
        $a+2b\le30$ (capital), $a\ge0$, $b\ge0$. Four half-planes, so a
        convex polygon.
  \item Vertices: $(0,0)$, $(20,0)$, $(0,15)$, and the intersection of
        $2a+b=40$ with $a+2b=30$. Solving: $b=40-2a$ and
        $a+80-4a=30$, so $a=\tfrac{50}{3}$ and $b=\tfrac{20}{3}$. A linear
        objective has parallel straight level lines; pushing them across a
        convex polygon, the last point of contact is a vertex (or a whole
        edge, if the level lines happen to be parallel to it). Hence
        checking the vertices suffices --- the fundamental theorem of
        linear programming.
  \item The two lines meet where $100-2q=20+q$, i.e. $q=\tfrac{80}{3}$ with
        $p=\tfrac{140}{3}$. The region is the triangle with vertices
        $(0,20)$, $(0,100)$ and
        $\bigl(\tfrac{80}{3},\tfrac{140}{3}\bigr)$, so the largest feasible
        quantity is $q=\tfrac{80}{3}\approx26.7$ --- beyond it no price
        satisfies both buyer and seller, which is exactly the competitive
        equilibrium quantity.
\end{solutions}

\section{Additional Problems}

\begin{problems}
  \item Sketch $3x-2y>6$ and say which side of the line contains the
        origin.
  \item Sketch $y\le -x+5$ together with $x\ge1$.
  \item Sketch $(x+1)^{2}+(y-2)^{2}<4$ and give centre and radius.
  \item Sketch $x^{2}+y^{2}\ge 1$ and $x^{2}+y^{2}\le 4$ together.
  \item Sketch $y\ge x^{2}-4$.
  \item Sketch $y\le 2^{x}$.
  \item Sketch $y\ge\ln x$ and state the domain restriction it forces.
  \item Sketch $xy\le 4$ on the first quadrant.
  \item Sketch $|x-1|\le 2$.
  \item Sketch $|x|+|y|\ge 1$ and decide convexity.
  \item Sketch $y>x^{2}$ and $y<x+6$, and find the $x$-range of the
        region.
  \item Sketch $x^{2}+4y^{2}\le 16$ and $y\ge 0$.
  \item Write inequalities for the closed triangle with vertices $(0,0)$,
        $(4,0)$, $(0,3)$.
  \item Write inequalities for the closed disc of radius $2$ centred at
        $(1,1)$, intersected with the first quadrant.
  \item Write inequalities for the region above $y=x^{2}$ and below
        $y=2x+3$, and give the $x$-range.
  \item Write inequalities for the horizontal strip between $y=-1$ and
        $y=3$, once with an absolute value and once without.
  \item Write inequalities for everything outside the ellipse
        $x^{2}/4+y^{2}=1$ but inside the circle $x^{2}+y^{2}=9$.
  \item Is the region $\{y\ge x^{2}\}\cap\{y\le 4\}$ bounded? Give its
        extreme $x$- and $y$-values.
  \item Give a system of three inequalities whose solution set is a single
        point.
  \item Give a system of three inequalities whose solution set is an
        unbounded convex region with exactly two vertices.
  \item Show that $\{(x,y): y\ge \mathrm{e}^{x}\}$ is convex.
  \item Show that $\{(x,y): y\ge \sqrt{x},\ x\ge0\}$ is \emph{not}
        convex, and say which nearby set \emph{is}.
  \item \textbf{Two constraints, two goods.} Sketch
        $x+y\le 10$, \;$2x+y\le 16$, \;$x,y\ge0$, and give all vertices.
  \item \textbf{Diet problem.} A feed must supply at least $8$ units of
        protein and $12$ of energy. Grain gives $2$ protein and $2$ energy
        per kg; soy gives $1$ protein and $4$ energy per kg. Write the
        feasible set, sketch it, and say why it is unbounded.
  \item \textbf{Minimum-wage wedge.} Sketch
        $\{(L,w): w\ge 10,\ w\le 30-L,\ L\ge0\}$ and interpret the region
        and its largest $L$.
\end{problems}

\section{Additional Solutions}

\begin{solutions}
  \item The open half-plane strictly below the line $y=\tfrac32x-3$.
        Testing $(0,0)$: $0>6$ is false, so the origin is \emph{not} in the
        region; it lies on the other side.
  \item The intersection of the half-plane on or below $y=5-x$ with the
        half-plane $x\ge1$: an unbounded wedge opening down-right, with
        corner $(1,4)$.
  \item The \emph{open} disc of radius $2$ centred at $(-1,2)$; the
        boundary circle is dashed and excluded.
  \item The closed annulus $1\le x^{2}+y^{2}\le4$ between the two circles.
        Not convex.
  \item The region on or above the parabola $y=x^{2}-4$, whose vertex is
        $(0,-4)$ and which crosses the axis at $x=\pm2$. Convex.
  \item Everything on or below the increasing curve $y=2^{x}$, which
        passes through $(0,1)$ and hugs the $x$-axis to the left. Since
        $2^{x}$ is \emph{convex}, the region below it is \textbf{not}
        convex: $(-3,\tfrac18)$ and $(3,8)$ both satisfy the inequality,
        but their midpoint $(0,4.0625)$ does not, since $2^{0}=1<4.0625$.
  \item The region on or above $y=\ln x$; the logarithm's domain confines
        it to the open right half-plane $x>0$. Since $\ln x$ is
        \emph{concave}, the region above it is \textbf{not} convex: the
        boundary points $(0.1,\ln 0.1)$ and $(10,\ln 10)$ have midpoint
        $(5.05,0)$, and $0<\ln 5.05\approx1.62$, so the midpoint is
        outside.
  \item The part of the first quadrant on or below the hyperbola $y=4/x$
        --- the sublevel set $\{xy\le4\}$. \textbf{Not} convex: $(0.5,8)$
        and $(8,0.5)$ both lie on the boundary, but their midpoint
        $(4.25,4.25)$ has $xy\approx18.1>4$. (Equivalently: $4/x$ is convex
        on $x>0$, so the region below it is not.)
  \item $-2\le x-1\le2$, i.e. the closed vertical strip $-1\le x\le3$.
        Convex.
  \item Everything on or outside the diamond $|x|+|y|=1$. Not convex ---
        the segment from $(1,0)$ to $(-1,0)$ passes through the origin,
        where $|x|+|y|=0<1$.
  \item Strictly between the parabola and the line; they meet where
        $x^{2}=x+6$, i.e. $x=-2$ and $x=3$, so the region spans
        $-2<x<3$. Both boundaries dashed.
  \item $\dfrac{x^{2}}{16}+\dfrac{y^{2}}{4}\le1$ with $y\ge0$: the upper
        half of an elliptical disc with semi-axes $4$ and $2$.
  \item $x\ge0$, \;$y\ge0$, \;$3x+4y\le12$ (the line through $(4,0)$ and
        $(0,3)$).
  \item $(x-1)^{2}+(y-1)^{2}\le4$, \;$x\ge0$, \;$y\ge0$.
  \item $y\ge x^{2}$ and $y\le 2x+3$; the boundaries meet where
        $x^{2}=2x+3$, i.e. $x=-1$ and $x=3$, so $-1\le x\le3$.
  \item With an absolute value: $|y-1|\le2$. Without: $y\ge-1$ and
        $y\le3$.
  \item $\dfrac{x^{2}}{4}+y^{2}\ge1$ \;and\; $x^{2}+y^{2}\le9$: the disc of
        radius $3$ with the elliptical disc punched out. Not convex.
  \item Yes, bounded: $y$ ranges over $[0,4]$ and $x$ over $[-2,2]$, since
        $y\ge x^{2}$ together with $y\le4$ forces $x^{2}\le4$. The extreme
        points are $(\pm2,4)$ and $(0,0)$.
  \item For example $x\ge0$, \;$y\ge0$, \;$x+y\le0$: the first two force
        $x,y\ge0$ and the third then forces $x=y=0$, so the solution set is
        exactly $\{(0,0)\}$.
  \item $x\ge0$, \;$y\ge0$, \;$x+y\ge2$. The boundary consists of the
        ray $\{x=0,\ y\ge2\}$, the segment from $(0,2)$ to $(2,0)$, and the
        ray $\{y=0,\ x\ge2\}$, so there are exactly two corners, $(0,2)$
        and $(2,0)$. The region is unbounded (it contains every large
        point of the quadrant) and convex, being an intersection of three
        half-planes.
  \item $\mathrm{e}^{x}$ is a convex function, and the set lying \emph{on
        or above} the graph of a convex function --- its epigraph --- is a
        convex set: if $y_{1}\ge\mathrm{e}^{x_{1}}$ and
        $y_{2}\ge\mathrm{e}^{x_{2}}$ then for $t\in[0,1]$
        \[
          ty_{1}+(1-t)y_{2}\ \ge\ t\mathrm{e}^{x_{1}}+(1-t)\mathrm{e}^{x_{2}}
          \ \ge\ \mathrm{e}^{tx_{1}+(1-t)x_{2}},
        \]
        the last step being exactly convexity of $\mathrm{e}^{x}$.
  \item $\sqrt{x}$ is \emph{concave}, so the region \emph{above} it is
        not convex. Concretely, $(0,0)$ and $(4,2)$ both lie on the
        boundary and so belong to the set, but their midpoint $(2,1)$ does
        not, since $\sqrt2\approx1.414>1$. The set that \emph{is} convex is
        the one below the curve, $\{y\le\sqrt{x},\ x\ge0\}$ --- the
        hypograph of a concave function.
  \item The lines meet where $x+y=10$ and $2x+y=16$, i.e. $x=6$, $y=4$.
        Vertices: $(0,0)$, $(8,0)$, $(6,4)$, $(0,10)$.
  \item With $g$ kg of grain and $s$ kg of soy: $2g+s\ge8$,
        $2g+4s\ge12$, $g\ge0$, $s\ge0$. It is unbounded because the
        constraints are all ``at least'' --- buying arbitrarily much feed is
        always still feasible, so the region extends indefinitely up and to
        the right. A cost-minimising objective is nevertheless attained,
        at one of the two vertices on the lower-left boundary.
  \item The region is the triangle bounded below by $w=10$, above by
        $w=30-L$, and on the left by $L=0$; the two sloping boundaries
        meet at $L=20$, $w=10$. It shows the employment levels compatible
        with a wage floor of $10$ when labour demand is $w=30-L$: the
        largest feasible $L$ is $20$, the employment the floor permits.
        Any $L$ beyond that would require a wage below the minimum.
\end{solutions}
